\documentclass[a4paper,12pt]{article}
\usepackage[T2A]{fontenc}
\usepackage[utf8]{inputenc}
\usepackage[russian]{babel}
\usepackage{amsmath,amssymb,mathtools}
\usepackage{helvet}
\renewcommand{\familydefault}{\sfdefault}
\usepackage{icomma}
\usepackage[left=20mm,right=20mm,top=23mm,bottom=20mm,headheight=25pt]{geometry}
\usepackage{microtype}
\usepackage{tikz}
\usetikzlibrary{arrows.meta,positioning,shapes.geometric,calc}
\usepackage{tabularx,booktabs}
\usepackage{enumitem}
\usepackage{parskip}
\usepackage{needspace}
\usepackage{fancyhdr}
\usepackage{setspace}
\usepackage{xcolor}
\usepackage{hyperref}
\hypersetup{hidelinks,pdftitle={Чек-лист. Текстовые задачи и алгебраические преобразования},pdfauthor={Лёвин Артём Александрович},pdfsubject={19 ноября 2026 — Ярослав Верещагин}}
\definecolor{accent}{HTML}{008080}
\definecolor{darkaccent}{HTML}{005C61}
\definecolor{textgray}{HTML}{333333}
\definecolor{rulegray}{HTML}{C8D2D2}
\color{textgray}
\onehalfspacing
\setlength{\abovedisplayskip}{6pt}
\setlength{\belowdisplayskip}{6pt}
\setlength{\abovedisplayshortskip}{4pt}
\setlength{\belowdisplayshortskip}{4pt}
\setlength{\parindent}{0pt}
\emergencystretch=1.5em
\setlist[itemize]{leftmargin=6mm,itemsep=1pt,topsep=2pt}
\setlist[enumerate]{leftmargin=8mm,itemsep=4pt,topsep=3pt}
\pagestyle{fancy}
\fancyhf{}
\fancyhead[L]{\small\color{darkaccent} Текстовые задачи и алгебра}
\fancyhead[R]{\small Ярослав Верещагин}
\renewcommand{\headrulewidth}{0.3pt}
\renewcommand{\headrule}{\hbox to\headwidth{\color{rulegray}\leaders\hrule height \headrulewidth\hfill}}
\fancyfoot[C]{\small\thepage}
\renewcommand{\footrulewidth}{0pt}
\newcommand{\checkitem}{\ensuremath{\square}\enspace}
\newcommand{\keyidea}[1]{\textcolor{darkaccent}{\bfseries #1}}
\newcommand{\sectionline}{\vspace{0.3em}\noindent\textcolor{accent}{\rule{\linewidth}{0.5pt}}\vspace{0.2em}}
\renewcommand{\section}[1]{\Needspace{5\baselineskip}\par\vspace{0.65em}\noindent{\Large\color{darkaccent}\bfseries #1}\par\vspace{0.15em}\sectionline}
\newcommand{\subhead}[1]{\par\vspace{0.4em}{\large\bfseries #1}\par\vspace{0.06em}}
\newcommand{\smallnote}[1]{{\small\color{textgray}#1}}
\newcommand{\qa}[2]{\noindent\textbf{#1} #2}
\begin{document}
\thispagestyle{empty}
\begin{center}
\vspace*{21mm}
{\large\color{darkaccent} МАТЕМАТИКА}\par
\vspace{25mm}
{\Huge\mdseries Чек-лист}\par
\vspace{17mm}
{\LARGE\color{accent}\begin{tabular}{c}
Текстовые задачи на движение\\
и совместную работу\\
Алгебраические преобразования
\end{tabular}}\par
\vspace{22mm}
{\large Подготовка к ЕГЭ по профильной математике}\par
\vfill
{\Large Ярослав Верещагин}\par
\vspace{12mm}
{\large 19 ноября 2026 г.}\par
\vspace{22mm}
{\normalsize Лёвин Артём Александрович\\[-0.1em]
эксклюзивно для Ярослава Верещагина}\par
\vspace{16mm}
\begin{tikzpicture}[x=1cm,y=1cm]
\draw[accent,line width=0.9pt] (-5.3,0) .. controls (-4.0,0.72) and (-2.2,-0.57) .. (-0.95,0)
.. controls (0.35,0.63) and (2.0,-0.57) .. (3.25,0)
.. controls (3.95,0.29) and (4.72,0.2) .. (5.3,0);
\end{tikzpicture}
\vspace*{3mm}
\end{center}
\clearpage

\section{1. Алгебра перед текстовой задачей}
\subhead{Степени и знаки}
\checkitem При возведении \emph{отрицательного числа} в чётную степень результат положительный, в нечётную — отрицательный:
\[
(-10)^4=10\,000,\qquad (-10)^3=-1000.
\]
Пример из занятия:
\[
0{,}6\cdot(-10)^4+4\cdot(-10)^3-90
=6000-4000-90=1910.
\]

\subhead{Две разные формулы, которые нельзя смешивать}
\checkitem \keyidea{Куб разности} — степень относится ко всей скобке:
\[
(a-b)^3=a^3-3a^2b+3ab^2-b^3.
\]
Например, \((2-x)^3=8-12x+6x^2-x^3\).

\checkitem \keyidea{Разность кубов} — два куба уже вычислены по отдельности:
\[
a^3-b^3=(a-b)(a^2+ab+b^2).
\]
Например, \(8-x^3=(2-x)(x^2+2x+4)\).

\keyidea{Ошибка:} запись \((a-b)^3=a^3-b^3\) неверна в общем случае.

\subhead{Как убрать куб из обеих частей уравнения}
\checkitem Если \(A^3=B^3\), то \(A=B\) для действительных чисел. В частности,
\[
 x^6=(8-7x)^3
 \ \Longleftrightarrow\ (x^2)^3=(8-7x)^3
 \ \Longleftrightarrow\ x^2=8-7x.
\]
Получаем \(x^2+7x-8=0\). Это обычное квадратное уравнение.

\subhead{Порядок решения квадратного уравнения}
\checkitem Приведите уравнение к виду \(ax^2+bx+c=0\), \(a\ne0\); \emph{явно выпишите} коэффициенты \(a,b,c\). Затем:
\[
D=b^2-4ac,\qquad x_{1,2}=\frac{-b\pm\sqrt{D}}{2a}\quad(D\ge0).
\]
Здесь \(a=1\), \(b=7\), \(c=-8\), поэтому
\[
D=49+32=81,\qquad
x_1=\frac{-7+9}{2}=1,\quad x_2=\frac{-7-9}{2}=-8.
\]
\checkitem Проверяйте корни подстановкой; в текстовых задачах дополнительно проверяйте \emph{смысл и ограничения}.

\section{2. Движение по реке}
\keyidea{Базовая формула:} расстояние равно скорости, умноженной на время:
\[
s=vt,\qquad t=\frac{s}{v},\qquad v=\frac{s}{t}.
\]

\checkitem Пусть \(x\) — собственная скорость лодки, а \(c\) — скорость течения. Тогда
\[
v_{\text{по теч.}}=x+c,\qquad v_{\text{против теч.}}=x-c,
\qquad x>c.
\]
Плот без собственного двигателя движется со скоростью течения: \(v_{\text{плота}}=c\).

\smallskip
\noindent\textbf{Таблица для задачи о движении:}\par
\smallskip
\renewcommand{\arraystretch}{1.35}
\begin{tabularx}{\linewidth}{@{}Xccc@{}}
\toprule
Участок пути & Скорость & Расстояние & Время \\
\midrule
По течению & \(x+c\) & \(s_1\) & \(\dfrac{s_1}{x+c}\)\\
Против течения & \(x-c\) & \(s_2\) & \(\dfrac{s_2}{x-c}\)\\
Плот & \(c\) & \(s_3\) & \(\dfrac{s_3}{c}\)\\
\bottomrule
\end{tabularx}
\renewcommand{\arraystretch}{1}

\subhead{Пример 1. Два участка пути}
Баржа прошла \(60\) км по течению и \(20\) км против течения за \(7\) часов. Течение — \(1\) км/ч. Найдите собственную скорость баржи.

Обозначим её через \(x>1\). Сумма времени на двух участках равна семи часам:
\[
\frac{60}{x+1}+\frac{20}{x-1}=7.
\]
Приводим к общему знаменателю \((x+1)(x-1)=x^2-1\):
\[
60(x-1)+20(x+1)=7(x^2-1),
\]
\[
7x^2-80x+33=0,\qquad D=5476=74^2.
\]
\[
x_1=11,\qquad x_2=\frac37.
\]
\keyidea{Ответ:} \(11\) км/ч. Число \(\tfrac37\) не подходит: скорость лодки в стоячей воде должна быть больше скорости течения.

\subhead{Проверка модели}
\checkitem На каждом участке используйте \emph{свою} скорость. \checkitem Если говорится «на весь путь», складывайте времена. \checkitem Минуты перед вычислениями переводите в часы.

\section{3. Плот вышел раньше лодки}
\subhead{Пример 2. Сначала сравниваем времена}
Лодка прошла \(108\) км по течению и столько же против течения. Плот, вышедший на \(2\) часа раньше, успел проплыть \(51\) км. Скорость течения — \(3\) км/ч. Найдите собственную скорость лодки.

\noindent\textbf{Шаг 1.} \(x>3\) — собственная скорость лодки. Время плота:
\[
t_{\text{плота}}=\frac{51}{3}=17\ \text{ч}.
\]
Плот вышел раньше, следовательно, он двигался на два часа \emph{дольше}. Время лодки равно \(17-2=15\) часам.

\noindent\textbf{Шаг 2.} Записываем путь, скорость и время для лодки:
\[
\frac{108}{x+3}+\frac{108}{x-3}=15.
\]
\noindent\textbf{Шаг 3.} Умножаем на \((x+3)(x-3)\). Скобки сохраняем:
\[
108(x-3)+108(x+3)=15(x^2-9).
\]
\[
216x=15x^2-135
\quad\Longleftrightarrow\quad
5x^2-72x-45=0.
\]
\noindent\textbf{Шаг 4.} Дискриминант и отбор корней:
\[
D=(-72)^2-4\cdot5\cdot(-45)=6084=78^2,
\]
\[
x_{1,2}=\frac{72\pm78}{10}
\quad\Longrightarrow\quad
x_1=15,\qquad x_2=-0{,}6.
\]
Второй корень не является допустимой скоростью. \keyidea{Ответ: \(15\) км/ч.}

\section{4. Не потерять смысл при составлении уравнения}
\begin{itemize}
\item \checkitem \textbf{Кто отправился раньше?} У него время движения больше, если события завершились одновременно.
\item \checkitem \textbf{Что одинаково?} Сумма времён, одно и то же расстояние, разница во времени — в зависимости от условия.
\item \checkitem \textbf{Где нужны скобки?} Например, \(-108(x-3)=-108x+324\), а не \(-108x-324\).
\item \checkitem \textbf{Какой физический смысл у \(x\)?} Скорость, время, расстояние и производительность не могут быть отрицательными.
\end{itemize}

\section{5. Задачи на совместную работу}
\keyidea{Три величины:} производительность \(p\), время \(t\), объём работы \(A\):
\[
A=pt,\qquad p=\frac At,\qquad t=\frac Ap.
\]
\checkitem Если всю работу принимают за единицу, то производительность исполнителя, который выполняет её за \(t\) часов, равна \(\tfrac1t\) работы в час.

\checkitem \keyidea{При одновременной независимой работе складываются производительности}, а не времена:
\[
p_{\text{вместе}}=p_1+p_2,
\qquad \frac1{t_{\text{вместе}}}=\frac1{t_1}+\frac1{t_2}.
\]

\subhead{Пример 3. Два мастера}
Один мастер выполняет заказ за \(3\) часа, второй — за \(6\) часов. Тогда
\[
\frac1x=\frac13+\frac16=\frac12
\quad\Longrightarrow\quad x=2\ \text{ч}.
\]
\keyidea{Проверка здравым смыслом:} вместе мастера должны закончить раньше, чем каждый по отдельности.

\subhead{Пример 4. Неизвестна скорость второго работника}
Два работника вместе выполняют задачу за \(8\) минут; первый один — за \(24\) минуты. Пусть второй один справляется за \(x\) минут:
\[
\frac18=\frac1{24}+\frac1x
\quad\Longrightarrow\quad
\frac1x=\frac18-\frac1{24}=\frac1{12}
\quad\Longrightarrow\quad x=12.
\]

\section{6. План и фактическая производительность}
\checkitem Если работа занимает \(A\) единиц, а производительность равна \(p\) единицам в день, то время равно \(\tfrac Ap\) дням.

\subhead{Пример 5. Набор текста}
Нужно набрать \(300\) страниц. Фактически наборщик печатает на \(5\) страниц в день больше плана и завершает работу на \(3\) дня раньше. Пусть \(x>0\) — плановая производительность:
\[
\frac{300}{x}-\frac{300}{x+5}=3.
\]
Меньшее фактическое время \(\tfrac{300}{x+5}\) вычитаем из большего планового:
\[
\frac{1500}{x(x+5)}=3
\ \Longleftrightarrow\ x^2+5x-500=0.
\]
Корни \(x=20\) и \(x=-25\); отрицательное значение не подходит.
\keyidea{Ответ:} \(20\) страниц в день по плану.

\section{7. Производительность труб и короткая самопроверка}
\subhead{Пример 6. Две трубы}
Первая труба пропускает на \(12\) литров в минуту меньше второй и наполняет бак объёмом \(160\) литров на \(12\) минут дольше. Пусть \(x>0\) — производительность первой трубы. Тогда у второй она равна \(x+12\):
\[
\frac{160}{x}-\frac{160}{x+12}=12.
\]
\[
\frac{1920}{x(x+12)}=12
\quad\Longleftrightarrow\quad
x^2+12x-160=0.
\]
\[
D=784=28^2,\qquad x_1=8,\quad x_2=-20.
\]
\keyidea{Ответ:} \(8\) л/мин. Вторая труба работает быстрее: \(20\) л/мин.

\subhead{Как оформлять такие задачи на экзамене}
\renewcommand{\arraystretch}{1.25}
\begin{tabularx}{\linewidth}{@{}lXXX@{}}
\toprule
Случай & Производительность & Время & Результат \\
\midrule
По плану & \(x\) & \(300/x\) & \(300\) страниц\\
Сверх плана & \(x+5\) & \(300/(x+5)\) & \(300\) страниц\\
\bottomrule
\end{tabularx}
\renewcommand{\arraystretch}{1}

\clearpage
\thispagestyle{empty}
\begin{center}
{\LARGE\bfseries\color{darkaccent} Алгоритм: как составить уравнение по условию}
\end{center}
\vspace{1.1em}
\begin{center}
\begin{tikzpicture}[
  >=Latex,
  action/.style={rectangle,rounded corners=2pt,draw=accent,line width=0.75pt,align=center,fill=white,text width=61mm,minimum height=13mm,inner sep=5pt,font=\small},
  side/.style={action,text width=51mm,minimum height=17mm},
  decision/.style={diamond,aspect=2.5,draw=accent,line width=0.75pt,align=center,inner sep=1pt,text width=43mm,minimum height=14mm,font=\small},
  endpoint/.style={ellipse,draw=darkaccent,line width=0.8pt,align=center,minimum width=53mm,minimum height=10mm,font=\small},
  line/.style={draw=accent,line width=0.8pt},
  arr/.style={-{Latex[length=2.5mm]},draw=accent,line width=0.85pt}
]
\node[endpoint] (start) at (0,0) {Прочитать задачу};
\node[action] (read) at (0,-1.8) {Выделить неизвестное,\ единицы измерения и ограничения};
\node[decision] (choice) at (0,-4.1) {Какой тип задачи?};
\node[side] (move) at (-3.85,-6.8) {\textbf{Движение}\\Скорость, время, расстояние\\\(t=s/v\)};
\node[side] (work) at (3.85,-6.8) {\textbf{Работа}\\Производительность, время, результат\\\(t=A/p\)};
\node[action] (table) at (0,-9.35) {Построить таблицу и выразить времена};
\node[action] (relation) at (0,-11.15) {Записать связь из условия:\ сумма времён, разность или\ сумма производительностей};
\node[action] (solve) at (0,-13.25) {Решить уравнение, сохраняя скобки};
\node[action] (test) at (0,-15.1) {Проверить ОДЗ, смысл корней\ и единицы ответа};
\node[endpoint] (end) at (0,-16.8) {Записать ответ};
\draw[arr] (start) -- (read);
\draw[arr] (read) -- (choice);
\draw[arr] (choice.west) -| node[pos=0.28,above,font=\scriptsize]{движение} (move.north);
\draw[arr] (choice.east) -| node[pos=0.28,above,font=\scriptsize]{работа} (work.north);
\coordinate (join) at (0,-8.5);
\draw[line] (move.south) |- (join);
\draw[line] (work.south) |- (join);
\draw[arr] (join) -- (table.north);
\draw[arr] (table) -- (relation);
\draw[arr] (relation) -- (solve);
\draw[arr] (solve) -- (test);
\draw[arr] (test) -- (end);
\end{tikzpicture}
\end{center}
\clearpage

\section{Тренировочный блок}
\textbf{Путь А.} Решите десять заданий по возрастанию сложности. Запишите промежуточные преобразования и ограничения там, где они нужны. В этом блоке решения не приведены.

\begin{enumerate}[leftmargin=8mm,itemsep=0.55em]
\item Вычислите: \(0{,}5\cdot(-10)^4+3\cdot(-10)^3-80\).
\item Раскройте скобки: \((3-x)^3\).
\item Разложите на множители: \(27-x^3\).
\item Решите уравнение: \(x^6=(5-4x)^3\).
\item Первый мастер выполняет заказ за \(4\) часа, второй — за \(12\) часов. За сколько часов они выполнят заказ вместе?
\item Два исполнителя делают работу за \(10\) минут. Первый по отдельности выполняет её за \(30\) минут. За сколько минут сделает эту работу второй?
\item Катер прошёл \(28\) км по течению и \(30\) км против течения за \(5\) часов. Найдите собственную скорость катера, если скорость течения \(2\) км/ч.
\item Наборщик должен набрать \(240\) страниц. Ежедневно печатая на \(4\) страницы больше плана, он заканчивает на \(2\) дня раньше. Сколько страниц в день было запланировано?
\item Первая труба пропускает на \(10\) л/мин меньше второй. Бак на \(180\) литров первая труба заполняет на \(9\) минут дольше второй. Найдите производительность первой трубы.
\item Плот проплыл \(24\) км со скоростью течения \(2\) км/ч. Лодка, вышедшая на \(2\) часа позже плота, прошла \(48\) км по течению и \(48\) км против течения и закончила путь одновременно с плотом. Найдите собственную скорость лодки.
\end{enumerate}
\vspace{0.4em}
\subhead{Перед сверкой ответов}
\begin{itemize}[itemsep=0pt,topsep=0pt]
\item \checkitem Указаны ограничения и единицы измерения.
\item \checkitem Скобки при умножении на общий знаменатель не потеряны.
\item \checkitem В совместной работе складывались производительности, а не времена.
\item \checkitem Полученные корни проверены по смыслу условия.
\end{itemize}

\clearpage
\section{Ответы к тренировочному блоку}
\vspace{0.3em}
\renewcommand{\arraystretch}{1.42}
\begin{tabularx}{\linewidth}{@{}cX@{}}
\toprule
№ & Ответ \\
\midrule
1 & \(1920\)\\
2 & \(27-27x+9x^2-x^3\)\\
3 & \((3-x)(x^2+3x+9)\)\\
4 & \(x=-5\) или \(x=1\)\\
5 & \(3\) часа\\
6 & \(15\) минут\\
7 & \(12\) км/ч\\
8 & \(20\) страниц в день\\
9 & \(10\) л/мин\\
10 & \(10\) км/ч\\
\bottomrule
\end{tabularx}
\renewcommand{\arraystretch}{1}

\vspace{0.9em}
\subhead{Короткая памятка по ограничениям}
Во всех задачах на движение против течения собственная скорость лодки больше скорости течения. В задачах на работу и производительность выбирайте только положительные значения времени и производительности. В квадратных уравнениях все найденные корни нужно проверить по исходному условию.

\vfill
\begin{center}
{\small\color{darkaccent} Лёвин Артём Александрович\\ эксклюзивно для Ярослава Верещагина}
\end{center}
\end{document}
